\section{The price of a sphere tensor}\label{sec:sphere}

How expensive can a tensor be when the dimensions of its legs are
prescribed? Prepare the tensor at a site \(o\), then send each leg to its
destination \(x_i\) along a wire. In the logarithmic price \(H\), a wire
costs its length times the logarithm of its dimension. In particular,
for every nonzero sphere tensor of Section~\ref{sec:fermions},
\begin{equation}\label{eq:sphere-transport}
 H(C(\mathbf h;\mathbf z);x_1,\ldots,x_k)
 \le \sum_{i=1}^k d(o,x_i)\log\dim V_{h_i}.
\end{equation}
The marked points \(z_i\) determine the tensor's entries. The sites
\(x_i\) determine its price. For legs of log dimension at most \(Na_i\),
the cheapest such center gives an upper bound of
\(N\inf_o\sum_i a_i d(o,x_i)\). Under the standing hypothesis
\(\omega_F<9/4\), sphere tensors attain this bound to leading order
(Corollary~\ref{cor:sphere-caps}).

We also need to know which energies make a sphere tensor expensive
at a given placement. We put a quadratic penalty on each leg width.
Squared width is proportional to the energy recorded by the order of
a collision. Theorem~\ref{thm:sphere} identifies the
optimizing widths with distances from a weighted barycenter, whose
weights are set by the penalties.

Fix an algebraically closed field \(F\), distinct \(z_1,\ldots,z_k\in F\),
and sites \(\mathbf x=(x_1,\ldots,x_k)\) in a complete CAT(0) spectral
space \(X\) over \(F\), with \(k\ge2\). Write \(H_{\mathbf x}\) for the
price at this placement and \(D\) for its scale. Let \(c=\pi/\sqrt3\)
be the constant from Lemma~\ref{lem:dimension}. For \(N\ge1\) an integer,
define the normalized width of energy \(h\) by \(w_N(h)=c\sqrt h/N\).
For positive penalties \(\lambda_1,\ldots,\lambda_k\), set
\begin{align}
 B_N&=\max_{\mathbf h}
 \left\{N^{-1}H_{\mathbf x}(C(\mathbf h;\mathbf z))-
                      \sum_i\lambda_iw_N(h_i)^2\right\},
 \label{eq:sphere-objective}\\
 E&=\min_{o\in X}\sum_i\frac{d(o,x_i)^2}{4\lambda_i}.
 \label{eq:sphere-energy}
\end{align}
Zero tensors may be discarded in the maximum. The unique minimizer \(o\)
in~\eqref{eq:sphere-energy} is the barycenter with weights proportional
to \(1/\lambda_i\). It lies in the closed convex hull of the \(x_i\).
Write \(r_i=d(o,x_i)\) and \(\rho_N=\sqrt{\log(2N)/N}\).

\begin{theorem}[Sphere prices]\label{thm:sphere}
Suppose \(\omega_F<9/4\). Then
\begin{equation}\label{eq:sphere-estimate}
       0\le E-B_N\le O(\rho_N).
\end{equation}
For every choice of maximizing energies \(\mathbf e\) at each \(N\),
\begin{equation}\label{eq:sphere-profile}
 w_N(e_i)=\frac{r_i}{2\lambda_i}+O(\sqrt{\rho_N}),
 \qquad
 \frac{H_{\mathbf x}(C(\mathbf e;\mathbf z))}{N}=2E+O(\sqrt{\rho_N}).
\end{equation}
The constants are uniform for fixed \(k\), bounded scales \(D\), and all
\(\lambda_i\) in a fixed compact subinterval of \((0,\infty)\).
\end{theorem}

For two sites at distance \(L\), the barycenter divides the joining
segment in the ratio \(\lambda_1:\lambda_2\). Both limiting widths
are therefore \(L/[2(\lambda_1+\lambda_2)]\). More generally, to
weight site \(i\) by a mass \(a_i>0\), choose
\(\lambda_i=1/(2a_i)\). Its width then approaches \(a_i r_i\), and the
normalized price approaches \(\sum_i a_i r_i^2\). Section~\ref{sec:incidence}
uses these individual widths to measure the effect of several
simultaneous collisions.

By Lemma~\ref{lem:dimension}, \(\log\dim V_h\le c\sqrt h\).
Completing squares in~\eqref{eq:sphere-transport} gives, for any
nonzero \(C=C(\mathbf h;\mathbf z)\),
\begin{equation}\label{eq:sphere-squares}
 \frac{H_{\mathbf x}(C)}{N}-\sum_i\lambda_i w_N(h_i)^2
 \le E-\sum_i\lambda_i
             \left(w_N(h_i)-\frac{r_i}{2\lambda_i}\right)^2 .
\end{equation}
The vacuum tensor, with all energies zero, has price zero. Thus the
maximum exists, \(0\le B_N\le E\), and the optimizing widths are bounded
uniformly in the range of Theorem~\ref{thm:sphere}, since the original
sites have diameter at most \(D\) and \(r_i\le D\). If all sites
coincide, zero energies attain the maximum and
\eqref{eq:sphere-squares} proves the theorem. For the remaining case,
the proof finds one extra site \(y_N\) satisfying
\[
 d(x_i,y_N)\le 2\lambda_i w_N(e_i)+O(\rho_N)\qquad(1\le i\le k)
\]
for every fixed maximizing tuple \(\mathbf e\). Squaring and summing
these distances gives a lower bound on the optimizer's penalty.
The change in its penalty from scale \(N\) to \(N+1\) then gives the
lower bound on \(B_N\).

\subsection{Pricing one additional leg}

The extension must work for all collisions at once. Synchronization
reduces that requirement to a price bound for products of sphere
tensors in which all additional legs have the same marked position
\(s\). A principal part at \(s=z_i\) then allows all the additional
legs to be closed at the single site \(x_i\).

\begin{lemma}\label{lem:sphere-added}
Suppose \(\omega_F<9/4\), fix \(R>0\) and put \(\mu=\max_i\lambda_i\).
For every \(N\ge2\) the base price has an extension of scale \(D\) to
one more site \(y_N\) such that
\begin{multline}\label{eq:sphere-aux}
 H_{\mathbf x,y_N}\bigl(C((\mathbf h,g);(\mathbf z,s))\bigr)\\
 \le N\left(B_N+\sum_j\lambda_j w_N(h_j)^2+\mu w_N(g)^2\right)
       + O_{k,R}(D\log(2N))
\end{multline}
simultaneously for all \(h_j,g\le RN^2\) and \(s\in F\setminus\{z_j\}\).
\end{lemma}

\begin{proof}
Let the \(u\)th factor in a product of \(m\) sphere tensors have energies
\(\mathbf h_u,g_u\) and additional position \(s\). A closure inserts
one joint state in the tensor product of the \(m\) additional argument
spaces. By Corollary~\ref{cor:sharedprice} for rational families, it
suffices to bound every such closure at the shared generic parameter.
A strict violation would persist on an open set of values of \(s\)
and of the closure coordinates, by closedness of price sublevels.
We may therefore fix any joint state over \(F\), keep it constant
as \(s\) varies, and bound the closed product at every regular value.

If all \(g_u=0\), the closed product is a scalar multiple of the base
product and~\eqref{eq:sphere-objective} applies.
Otherwise a factor with \(g_u>0\) must contract at least one generator
at \(s\) with a generator at a fixed point in every matching. Those
contractions vanish at infinity. Thus the closed product is a rational
tensor vanishing at infinity, with poles only at the \(z_i\), and equals
the sum of its principal parts.

At the pole \(z_i\), fuse the additional point in each factor with its
\(i\)th point, using offsets \(1\) and \(0\), respectively. A term with
intermediate energies \(e_u\) has Laurent order \(\sum_u(e_u-h_{ui}-g_u)\), by
Proposition~\ref{prop:fusion}. After inserting the joint state, this term
is a restriction of a product of base sphere tensors in which the
\(i\)th energies have been replaced by \(e_u\). Indeed all \(m\) fusion
maps and the joint state now act at the single site \(x_i\). The price
of a term of negative order is therefore at most
\begin{align*}
 &\sum_u N\left(B_N+\sum_j\lambda_j w_N(h_{uj})^2
             +\lambda_i\bigl(w_N(e_u)^2-w_N(h_{ui})^2\bigr)\right)
 \\
 &\hspace{40pt}\le \sum_u N\left(B_N+\sum_j\lambda_j w_N(h_{uj})^2
                      +\mu w_N(g_u)^2\right).
\end{align*}
Indeed a principal-part term has \(\sum_u e_u<\sum_u(h_{ui}+g_u)\).
The identity \(w_N(h)^2=c^2h/N^2\) converts this single inequality
on total energy to the displayed bound, including factors whose
individual Laurent orders are positive. The integers \(2e_u\) have sum at most
\(4RmN^2\), giving at most
\(\binom{m+\lfloor4RmN^2\rfloor}{m}\) possibilities for each \(i\).
Thus the ordinary-sum price inequality costs \(O_{k,R}(mD\log(2N))\)
for all principal parts. This proves the required product test with
the error in~\eqref{eq:sphere-aux} per factor.
\end{proof}

\subsection{Recovering a wire by a collision}

\begin{proof}[Proof of Theorem~\ref{thm:sphere}]
Fix maximizing energies \(\mathbf e\), abbreviate \(u_i=w_N(e_i)\),
and write \(Q_N=\sum_i\lambda_i u_i^2\). Choose a half-integer
\(g\) nearest to \(N\log(2N)\). Lemma~\ref{lem:dimension} gives
\begin{equation}\label{eq:sphere-testwidth}
 v=w_N(g)\asymp\sqrt{\frac{\log(2N)}{N}},
 \qquad
 \log\dim V_g=Nv-O(\log(2N))
\end{equation}
for sufficiently large \(N\). We use this small but growing space to
probe the distance from \(y_N\) to each \(x_i\). The choice
\(v\asymp\rho_N\) balances the quadratic cost \(v^2\) of the added width
with the normalized interpolation cost \(O(\log(2N)/N)\).

Hold \(h_j=e_j\) for \(j\ne i\). Lemma~\ref{lem:width} says that the
fusion maps at \(z_i\) give a surjection
\begin{equation}\label{eq:sphere-surjection}
 \bigoplus_h V_h\longrightarrow\operatorname{Hom}(V_g,V_{e_i}),
 \qquad (a_h)_h\longmapsto\sum_h I_{e_i}(a_h,-),
\end{equation}
using only \(h\) with \(\sqrt h\le\sqrt{e_i}+\sqrt g+O(1)\).
Choose \(R\) in Lemma~\ref{lem:sphere-added} large enough for these
energies and all other \(e_j\). The bound on optimizing widths makes
one \(R\) sufficient throughout the uniform range of the theorem.
Choose the site \(y_N\) from that lemma. It works simultaneously
for every choice of \(i,h\) and every regular \(s\). In the single
sphere tensor with \(h_i=h\), extract the coefficient of
\((s-z_i)^{e_i-h-g}\) whenever the parities agree. Fusion identifies
this coefficient, up to a scalar sign, with \(C(\mathbf e;\mathbf z)\)
precomposed at \(i\) by \(I_{e_i}:V_h\otimes V_g\to V_{e_i}\).

The required order \(e_i-h-g\) can be positive. Each coordinate
function is bounded at infinity, and fusion bounds its pole
orders at the \(z_j\) by \(2RN^2\). A common denominator of degree
\(O_{k,R}(N^2)\) consequently turns every entry into a polynomial of
at most that degree. Interpolating at distinct regular points of \(F\)
recovers any Laurent coefficient as the same linear combination of
the corresponding tensor evaluations.

Identify \(a\otimes f\in V_{e_i}\otimes V_g^*\) with the map
\(b\mapsto f(b)a\). A right inverse in~\eqref{eq:sphere-surjection}
writes this map as a sum of fusion maps, linearly in \(a\otimes f\).
Precompose the recovered coefficients with the components of this
inverse at \(x_i\) and add. With \(b\in V_g\) at \(y_N\), the resulting
form is
\[
 (a\otimes f,b)\longmapsto f(b)\,
             C(\mathbf e;\mathbf z)(\ldots,a,\ldots).
\]
Figure~\ref{fig:sphere-wire} illustrates the resulting factorization.
The tensor is the optimizer times a wire of dimension \(\dim V_g\)
between \(x_i\) and \(y_N\). Its price is the sum of their logarithmic
prices.
All \(h\) and interpolation values together contribute only polynomially
many summands, so the ordinary-sum inequality costs \(O(\log(2N))\).
By \eqref{eq:sphere-aux}, producing the wire can increase the old
\(i\)th width from \(u_i\) to at most \(u_i+v+O(N^{-1})\) and adds the
new width \(v\). Subtract the optimizer's price \(N(B_N+Q_N)\) to get
\begin{align*}
 d(x_i,y_N)\frac{\log\dim V_g}{N}
 &\le \lambda_i\bigl((u_i+v+O(N^{-1}))^2-u_i^2\bigr)
           +\mu v^2+O\!\left(\frac{\log(2N)}{N}\right)\\
 &\le 2\lambda_i u_i v+O\!\left(v^2+\frac{\log(2N)}{N}\right).
\end{align*}

\begin{figure}[tbp]
\centering
\begin{tikzpicture}[font=\small,x=1cm,y=1cm,
 dot/.style={circle,fill=ink,inner sep=2.5pt},
 gate/.style={draw=ink,rounded corners=2pt,fill=ink!5,inner sep=3pt},
 wire/.style={line width=.9pt,draw=sea}]
 \node[font=\small\bfseries] at (-3.25,1.45) {Collision at a marked point};
 \node[dot,label=below:{$z_i$},label=above:{$V_h$}] (old) at (-4.45,0) {};
 \node[dot,label=below:{$s$},label=above:{$V_g$}] (new) at (-2.10,0) {};
 \draw[-{Stealth[length=5pt]},draw=rust,line width=.8pt]
       ([xshift=-3pt]new.west) -- ([xshift=3pt]old.east);
 \node[font=\footnotesize] at (-3.27,-.72)
       {coefficient gives \(I_{e_i}\colon V_h\otimes V_g\to V_{e_i}\)};

 \node[font=\small\bfseries] at (3.1,1.45) {Tensor at the spectral sites};
 \node[dot,label=below:{$a\otimes f$ at $x_i$}] (xi) at (2.73,-.15) {};
 \node[dot,label=below:{$b$ at $y_N$}] (y) at (4.77,-.15) {};
 \node[gate] (c) at (1.3,.58) {\(C(\mathbf e;\mathbf z)\)};
 \draw[draw=ink,line width=.8pt] (c) -- node[above right,inner sep=1pt] {\(a\)} (xi);
 \draw[wire] (xi) -- node[above] {\(f(b)\)} (y);
 \draw[draw=ink] (c.west) -- ++(-.28,.15);
 \draw[draw=ink] (c.west) -- ++(-.28,-.15);
\end{tikzpicture}
\caption{The collision at \(s=z_i\) supplies a fusion map for each
admissible \(h\). Their span contains \(b\mapsto f(b)a\).
Inserting this map gives the original tensor with input \(a\) and a
wire joining the two spectral sites. The other legs keep their
placements.}\label{fig:sphere-wire}
\end{figure}

The extended scale bounds \(d(x_i,y_N)\le D\). Hence
\eqref{eq:sphere-testwidth} yields \(d(x_i,y_N)\le2\lambda_i u_i+O(v)\).
All constants here are uniform as stated.

The extension can be realized in a complete CAT(0) spectral space.
The closed convex hull of the old sites is canonically isometric to
their original hull by Lemma~\ref{lem:geodesic}.
Projection onto it weakly decreases distances to the old sites.
Consequently
\begin{equation}\label{eq:sphere-Q}
 E\le\sum_i\frac{d(x_i,y_N)^2}{4\lambda_i}
     \le Q_N+O\!\left(\sqrt{\frac{\log(2N)}{N}}\right).
\end{equation}
This is a lower bound on the penalty paid by every optimizer. To turn
it into a lower bound on the optimized price, reuse the same energies
in the objective defining \(B_{N+1}\). Their price is unchanged, while
their penalty in \((N+1)B_{N+1}\) is \(N^2Q_N/(N+1)\), compared with
\(NQ_N\) in \(NB_N\). Thus
\[
 (N+1)B_{N+1}-NB_N\ge\frac{N}{N+1}Q_N.
\]
Sum this inequality and~\eqref{eq:sphere-Q} up to scale \(N\).
The errors sum to \(O(\sqrt{N\log(2N)})\) and \(E\) is uniformly
bounded. Together with \(B_N\le E\) this proves~\eqref{eq:sphere-estimate}.
Equation~\eqref{eq:sphere-squares} now gives the width estimates
in~\eqref{eq:sphere-profile}. The widths and their targets are bounded,
so \(Q_N=E+O(\sqrt{\rho_N})\). Adding \(B_N\) gives the asserted price.
\end{proof}

\subsection{Prescribed widths and uniform comparison}

\begin{corollary}[Prescribed widths]\label{cor:sphere-caps}
Under the hypotheses of Theorem~\ref{thm:sphere}, fix positive width
caps \(a_i\). Then
\begin{equation}\label{eq:sphere-caps}
 \lim_{N\to\infty}
 \frac{1}{N}\max_{\mathbf h:\,w_N(h_i)\le a_i}
            H_{\mathbf x}(C(\mathbf h;\mathbf z))
 = \inf_{o\in X}\sum_i a_i d(o,x_i).
\end{equation}
Convergence is uniform over all placements of bounded scale.
\end{corollary}

\begin{proof}
Transport gives the upper bound. Fix \(\tau>0\) and minimize
\(\sum_i a_i\sqrt{d(o,x_i)^2+\tau^2}\) on the closed convex hull of
the sites. For \(r\) between zero and the hull diameter the function \(r\mapsto\sqrt{r^2+\tau^2}-\epsilon r^2\) is
convex and increasing for sufficiently small \(\epsilon>0\). The CAT(0)
midpoint inequality then makes every minimizing sequence Cauchy, giving
a minimizer by completeness.
Metric projection shows that this is also a minimizer on \(X\).
Write \(r_i=d(o,x_i)\) at this minimizer and set
\(2\lambda_i=\sqrt{r_i^2+\tau^2}/a_i\). The tangent upper bound for the
concave function \(s\mapsto\sqrt{s+\tau^2}\) shows that \(o\) minimizes
\(\sum_i d(o,x_i)^2/(4\lambda_i)\) as well. Indeed the tangent
majorant for the smoothed objective agrees with it at \(o\) and differs
from this quadratic objective by a constant.

At scale bounded by \(D_0\), we have \(r_i\le D_0\). For fixed \(\tau\)
these penalties stay in a compact positive interval. The limiting
widths in~\eqref{eq:sphere-profile} are
\(a_i r_i/\sqrt{r_i^2+\tau^2}\), bounded strictly below \(a_i\) uniformly.
Thus every maximizing tuple at large \(N\) meets the caps. Its normalized
price is
\[
 2E+o(1)
 =\sum_i a_i\frac{r_i^2}{\sqrt{r_i^2+\tau^2}}+o(1)
 \ \ge\ \inf_{z\in X}\sum_i a_i d(z,x_i)-\tau\sum_i a_i+o(1).
\]
Let \(\tau\downarrow0\) after taking \(N\to\infty\).
\end{proof}

A maximizing energy tuple in Corollary~\ref{cor:sphere-caps} can depend
on the placement. The local construction of Section~\ref{subsec:plant-private}
compares a single tensor against all spectral prices simultaneously.
Taking a direct sum over the admissible energies gives that tensor at
a logarithmic cost in each leg dimension.

\begin{corollary}[A common comparison tensor]\label{cor:universal}
Suppose \(\omega_F<9/4\). Fix the \(k\) marked positions and positive
caps \(a_1,\ldots,a_k\). Let \(A_N\) be the direct sum of all nonzero
\(C(\mathbf h;\mathbf z)\) with \(w_N(h_i)\le a_i\) at every leg.
Then
\[
 \frac{H_{\mathbf x}(A_N)}{N}=\inf_{o\in X}\sum_i a_i d(o,x_i)+o(1)
\]
uniformly over placements of bounded scale in complete CAT(0)
spectral spaces over \(F\). The \(i\)th leg of \(A_N\) has log dimension
at most \(Na_i+O_{k,\mathbf a}(\log(2N))\).
\end{corollary}

\begin{proof}
Each leg admits \(O_{\mathbf a}(N^2)\) half-integer energies, so
there are \(O_{k,\mathbf a}(N^{2k})\) summands. Each summand is a
restriction of \(A_N\), and the direct-sum inequality bounds the price
of \(A_N\) by the largest summand price plus
\(O_{k,\mathbf a}(D\log(2N))\). Corollary~\ref{cor:sphere-caps} now
gives the stated limit. Each \(i\)th summand leg has dimension at most
\(e^{Na_i}\) by Lemma~\ref{lem:dimension}, so summing those dimensions
gives the claimed bound.
\end{proof}
